\documentclass{article}

\usepackage[margin=1in]{geometry}
\usepackage{booktabs}
\usepackage{array}
\usepackage{style/tablemarks}

\begin{document}

\section*{Tablemarks Usage Example}

This file demonstrates how to mark the best and second-best values in result
tables. Compile it from the project root with:

\begin{verbatim}
latexmk -pdf -outdir=output style/tablemarks.tex
\end{verbatim}

\subsection*{Column-wise Metrics}

\tmsetup{style=both}

\begin{table}[h]
  \centering
  \caption{Automatic column-wise marking.}
  \begin{tabular}{lccc}
    \toprule
    Method & PSNR$\uparrow$ & MS-SSIM$\uparrow$ & LPIPS$\downarrow$ \\
    \midrule
    D-NeRF & \tmcell{col-psnr}{38.18} & \tmcell{col-ssim}{0.9910} & \tmlowcell{col-lpips}{0.0120} \\
    TiNeuVox-B & \tmcell{col-psnr}{40.62} & \tmcell{col-ssim}{0.9968} & \tmlowcell{col-lpips}{0.0083} \\
    SC-GS & \tmcell{col-psnr}{44.91} & \tmcell{col-ssim}{0.9980} & \tmlowcell{col-lpips}{0.0166} \\
    Ours & \tmcell{col-psnr}{41.72} & \tmcell{col-ssim}{0.9970} & \tmlowcell{col-lpips}{0.0097} \\
    \bottomrule
  \end{tabular}
\end{table}

\subsection*{Custom Best and Second-best Styles}

Use \verb|best-style| and \verb|second-style| when the best and second-best
cells should use different visual treatments.

\tmsetup{
  best-style=both,
  second-style=background,
  best-color=blue!20,
  second-color=green!20,
  bold-line-width=0.30pt
}

\begin{table}[h]
  \centering
  \caption{Custom style: blue bold best, green second-best.}
  \begin{tabular}{lc}
    \toprule
    Method & Accuracy$\uparrow$ \\
    \midrule
    Method A & \tmcell{custom-acc}{82.1} \\
    Method B & \tmcell{custom-acc}{84.6} \\
    Method C & \tmcell{custom-acc}{83.7} \\
    Ours & \tmcell{custom-acc}{86.2} \\
    \bottomrule
  \end{tabular}
\end{table}

\subsection*{Heatmap Cells}

Use \verb|\tmheatcell| when higher values should have darker backgrounds, and
\verb|\tmlowheatcell| when lower values should have darker backgrounds. The
package collects each group's minimum and maximum values automatically.

\tmsetup{
  heat-color=blue,
  heat-min-intensity=8,
  heat-max-intensity=55
}

\begin{table}[h]
  \centering
  \caption{Heatmap backgrounds based on cell values.}
  \begin{tabular}{lcc}
    \toprule
    Method & Accuracy$\uparrow$ & Error$\downarrow$ \\
    \midrule
    Method A & \tmheatcell{heat-acc}{82.1} & \tmlowheatcell{heat-err}{17.9} \\
    Method B & \tmheatcell{heat-acc}{84.6} & \tmlowheatcell{heat-err}{15.4} \\
    Method C & \tmheatcell{heat-acc}{83.7} & \tmlowheatcell{heat-err}{16.3} \\
    Ours & \tmheatcell{heat-acc}{86.2} & \tmlowheatcell{heat-err}{13.8} \\
    \bottomrule
  \end{tabular}
\end{table}

\subsection*{Formatted Cells}

If the cell contains units, uncertainty, or extra formatting, \verb|\tmcell| uses the
first number in the displayed content as the comparison value. The package
collects these values from the table cells automatically, so there is no
separate value list to maintain.

\tmsetup{style=text}

\begin{table}[h]
  \centering
  \caption{Formatted cells with automatic numeric extraction.}
  \begin{tabular}{lc}
    \toprule
    Method & PSNR$\uparrow$ \\
    \midrule
    Baseline & \tmcell{fmt-psnr}{$28.47 \pm 0.15$} \\
    Method A & \tmcell{fmt-psnr}{$30.68 \pm 0.09$} \\
    Method B & \tmcell{fmt-psnr}{$31.40 \pm 0.11$} \\
    Ours & \tmcell{fmt-psnr}{$32.45 \pm 0.07$} \\
    \bottomrule
  \end{tabular}
\end{table}

\subsection*{Row-wise Metrics}

For row-wise marking, declare one group per row and reuse that group across the
cells in the row.

\tmsetup{style=background}

\begin{table}[h]
  \centering
  \caption{Automatic row-wise marking.}
  \begin{tabular}{lcccc}
    \toprule
    Scene & D-NeRF & TiNeuVox-B & SC-GS & Ours \\
    \midrule
    Scene A
      & \tmcell{row-scene-a}{31.02}
      & \tmcell{row-scene-a}{32.78}
      & \tmcell{row-scene-a}{35.51}
      & \tmcell{row-scene-a}{34.80} \\
    Scene B
      & \tmcell{row-scene-b}{29.65}
      & \tmcell{row-scene-b}{30.12}
      & \tmcell{row-scene-b}{31.24}
      & \tmcell{row-scene-b}{33.08} \\
    \bottomrule
  \end{tabular}
\end{table}

\end{document}
